% Canonical standalone description of the computational analyses.
% Consolidates and supersedes the five focused notes named in Section 1.
\documentclass[11pt]{article}
\usepackage[margin=1in]{geometry}
\usepackage{amsmath,amssymb}
\usepackage{booktabs}
\usepackage{array}
\usepackage{longtable}
\usepackage[colorlinks=true,linkcolor=blue,citecolor=blue,urlcolor=blue]{hyperref}

\title{Consolidated Computational Analysis:\
Dual-Compartment Bacteremia Model, Monte Carlo Procedure,\
and Deterministic Switch Localization}
\author{}
\date{4 September 2026}

\newcommand{\EC}{EC_{50,L}}
\newcommand{\EM}{E_{\max,L}}
\newcommand{\KI}{k_{\mathrm{immune}}}
\newcommand{\LOD}{L_{\mathrm{OD}}}

\begin{document}
\maketitle
\tableofcontents

\section{Purpose and scope}

This document is the canonical, standalone explanation of analyses previously
divided among five focused notes, which were removed after consolidation:
\begin{itemize}
\item \path{EC50_lzd_analysis.tex};
\item \path{ec50_threshold_comparison_summary.tex};
\item \path{MC_immune_response_procedure.tex};
\item \path{MC_scripts_comparison.tex}; and
\item \path{vancomycin_sensitive_clearance.tex}.
\end{itemize}
Obsolete descriptions from those notes are not carried forward. In
particular, the current model uses
\texttt{solve\_ivp}, not \texttt{odeint}; uses the current centralized growth
rates; and treats the limit of detection as a hard extinction/reactivation
boundary.

Three one-parameter Monte Carlo analyses are included:
\begin{enumerate}
\item linezolid half-maximal concentration, $\EC$;
\item linezolid maximum killing effect, $\EM$; and
\item host immune-clearance rate, $\KI$.
\end{enumerate}
Each analysis varies only the named parameter. It is therefore a univariate
sensitivity/variability analysis, not a joint population model in which all
parameters vary simultaneously. Initial conditions are fixed and are not an
analysis variable.

\section{Model and treatment protocol}

\subsection{State variables}

The state vector is
\begin{equation}
\mathbf{Y}(t)=\bigl[S_b(t),R_b(t),S_{\mathrm{res}}(t),R_{\mathrm{res}}(t)\bigr],
\end{equation}
where $S$ and $R$ denote vancomycin-sensitive and vancomycin-resistant
bacteria, and the subscripts denote blood and the bone reservoir. The fixed
starting state is
\begin{equation}
\mathbf{Y}(0)=[0,0,100,100]\ \mathrm{CFU/mL}.
\end{equation}
This starting state is required to run the model but is not sampled in the
present analyses.

The simulation lasts 1944~h (81~days): 21~days before treatment, 4~days of
vancomycin, 42~days of linezolid, and 14~days of post-linezolid follow-up.

\subsection{Drug-effect functions}

With Hill coefficients equal to one, the blood-compartment effects are
\begin{align}
E_V(t)&=\frac{E_{\max,V}V(t)}{EC_{50,V}+V(t)},\\
E_L(t)&=\frac{\EM L(t)}{\EC+L(t)}.
\end{align}
Larger $\EC$ means weaker linezolid potency at the same concentration,
whereas larger $\EM$ means a larger maximum killing capacity. Reservoir drug
concentrations are the corresponding blood concentrations multiplied by the
centralized tissue-penetration fractions. In the reservoir the linezolid
ceiling is further reduced by an independent efficacy factor,
$E_{\max,L}^{\mathrm{res}}=\alpha\,\EM$ with $\alpha=0.41$, representing
reduced activity against slow-growing biofilm and small-colony-variant
bacteria in bone; $\alpha$ is calibrated and is not derived from the growth
rates.

\subsection{Centralized configuration}

All scripts obtain the fixed configuration from \texttt{parameters.py}.
Values are shown to two significant figures.

\begin{table}[htbp]
\centering
\small
\begin{tabular}{@{}lll@{}}
\toprule
Quantity & Value & Interpretation \\
\midrule
$\rho_S$ & $1.1~\mathrm{h}^{-1}$ & sensitive growth, blood \\
$\rho_R$ & $0.99~\mathrm{h}^{-1}$ & resistant growth, blood \\
$\rho_{\mathrm{res},S}$ & $0.50~\mathrm{h}^{-1}$ & sensitive growth, reservoir \\
$\rho_{\mathrm{res},R}$ & $0.45~\mathrm{h}^{-1}$ & resistant growth, reservoir \\
$E_{\max,V}$ & $1.7~\mathrm{h}^{-1}$ & maximum vancomycin effect \\
$EC_{50,V}$ & $0.25~\mathrm{mg/L}$ & vancomycin half-maximal concentration \\
$\EM$ & $1.65~\mathrm{h}^{-1}$ & linezolid maximum effect (Schmidt et al.) \\
$\EC$ & $1.13~\mathrm{mg/L}$ & linezolid half-maximal concentration (Schmidt et al.) \\
$B_{\max,b}$ & $6.0\times10^3~\mathrm{CFU/mL}$ & blood carrying capacity \\
$B_{\max,\mathrm{res}}$ & $1.0\times10^7~\mathrm{CFU/mL}$ & reservoir carrying capacity \\
$\phi_V$ & $0.15$ & reservoir vancomycin penetration \\
$\phi_L$ & $0.40$ & reservoir linezolid penetration \\
$\alpha$ & $0.41$ & reservoir linezolid efficacy factor \\
$k_{r\to b}$ & $1.0\times10^{-7}~\mathrm{h}^{-1}$ & reservoir-to-blood transfer \\
$\KI$ & $0.10~\mathrm{h}^{-1}$ & baseline immune clearance \\
$\epsilon_{\mathrm{res}}$ & $0.50$ & reservoir immune multiplier \\
$\LOD$ & $10~\mathrm{CFU/mL}$ & hard detection boundary \\
\bottomrule
\end{tabular}
\caption{Current shared configuration. The resistant-strain fitness cost is
10\% in both blood and the reservoir. Exchange is one-way, reservoir to
blood; there is no blood-to-reservoir term.}
\label{tab:parameters}
\end{table}

\section{Vancomycin-sensitive-strain clearance}

Vancomycin acts directly on $S_b$ and $S_{\mathrm{res}}$ but not on the
resistant populations. At low blood burden and saturating vancomycin, the
approximate per-capita sensitive-blood rate is
\begin{equation}
\rho_S-\KI-E_{\max,V}=1.1-0.10-1.7=-0.70~\mathrm{h}^{-1}.
\end{equation}
The negative value explains the rapid decline of the sensitive strain after
vancomycin begins. Removal of sensitive bacteria also releases shared carrying
capacity, allowing resistant bacteria to expand during vancomycin treatment.

The current baseline hybrid simulation gives the following values. A zero is
an exact model state produced by the hard-LOD rule, not an arbitrarily small
positive value rounded for display.

\begin{table}[htbp]
\centering
\small
\begin{tabular}{@{}lrrrr@{}}
\toprule
Time & $S_b$ & $R_b$ & $S_{\mathrm{res}}$ & $R_{\mathrm{res}}$ \\
\midrule
Vancomycin start & $4.7\times10^3$ & $0$ & $7.8\times10^6$ & $2.3\times10^5$ \\
$+24$ h & $0$ & $5.6\times10^2$ & $0$ & $8.9\times10^6$ \\
$+48$ h & $0$ & $5.4\times10^3$ & $0$ & $8.9\times10^6$ \\
$+72$ h & $0$ & $5.4\times10^3$ & $0$ & $8.9\times10^6$ \\
Vancomycin end & $0$ & $4.8\times10^3$ & $0$ & $8.6\times10^6$ \\
\bottomrule
\end{tabular}
\caption{Current baseline burdens during the four-day vancomycin course,
in CFU/mL.}
\label{tab:vancomycin-course}
\end{table}

Thus vancomycin clears both sensitive compartments within the first day of
treatment under the current calibration. Resistant blood bacteria are
undetectable before treatment and become detectable once vancomycin frees
the shared carrying capacity. Subsequent linezolid treatment is
responsible for control of the resistant populations.

\section{Numerical integration and the hard LOD}

All current analyses call \texttt{solver.integrate\_model}, which wraps
\texttt{scipy.integrate.solve\_ivp}. The configured method is LSODA, with
relative tolerance $10^{-6}$, absolute tolerance $10^{-8}$ CFU/mL, and
maximum internal step 4.0~h. The canonical output grid starts from 1450
evenly spaced values and inserts the three treatment boundaries, producing
1453 output times.

The solver is hybrid because it combines continuous ODE integration with
discrete state changes:
\begin{enumerate}
\item While a population is detectable, it follows the continuous ODE.
\item A downward event at $\LOD=10$ sets that population to exactly zero.
\item At zero, the population has no autonomous growth, drug killing, immune
clearance, or outgoing transfer.
\item Positive incoming transfer from a detectable paired compartment is
allowed to accumulate.
\item When accumulated arrivals reach the LOD, an upward event reactivates
the population and continuous dynamics resume.
\end{enumerate}
This construction prevents a fractional sub-LOD remnant from manufacturing a
relapse while still permitting a persistent reservoir to reseed blood.

\section{Two related but distinct analyses}

The deterministic and Monte Carlo calculations answer different questions.

\begin{description}
\item[Deterministic switch analysis.] If every value except one parameter is
fixed, where does the model change between final reservoir clearance and
final reservoir persistence?
\item[Monte Carlo analysis.] If that parameter differs among 1000 virtual
patients according to a specified distribution, what fraction experiences
bacteremia relapse, persistent bacteremia, or resolution?
\end{description}

The deterministic switch helps explain the Monte Carlo percentage, but it
does not assign the clinical outcome. Clinical outcomes are calculated from
each patient's complete blood trajectory.

\section{Deterministic switch localization with \texttt{brentq}}

\subsection{The quantity being localized}

Let $\theta$ be the parameter being varied and let $t_f$ be the end of the
81-day simulation. Define
\begin{equation}
F(\theta)=R_{\mathrm{res}}(t_f;\theta)-\LOD.
\label{eq:brent-objective}
\end{equation}
The sign is interpreted as
\begin{align}
F(\theta)<0 &\quad\Longrightarrow\quad \text{reservoir cleared},\\
F(\theta)>0 &\quad\Longrightarrow\quad \text{reservoir detectable}.
\end{align}
Because the hybrid solver returns zero rather than a positive sub-LOD value,
the clearance side generally has $F=-10$ rather than a value approaching
zero smoothly.

\subsection{What the sign does---and does not---mean}
\label{sec:sign-meaning}

The function $F$ contains the \emph{reservoir} burden at the final observation,
not the blood burden after linezolid. Its sign must therefore be read only as
follows:
\begin{center}
\fbox{\begin{minipage}{0.88\textwidth}
\begin{align*}
F(\theta)<0
&\quad\Longrightarrow\quad
R_{\mathrm{res}}(t_f)<10
\quad\text{(reservoir cleared at day 81)},\\
F(\theta)>0
&\quad\Longrightarrow\quad
R_{\mathrm{res}}(t_f)>10
\quad\text{(reservoir persists at day 81)}.
\end{align*}
Neither statement, by itself, says whether bacteremia occurred after
linezolid. In particular, a negative value of $F$ does \textbf{not} mean that
post-linezolid bacteremia occurred. A positive value of $F$ also does
\textbf{not} automatically mean relapse.
\end{minipage}}
\end{center}

Two times and two compartments are involved:
\begin{center}
\begin{tabular}{@{}lll@{}}
\toprule
Quantity & Time & Question answered \\
\midrule
$R_{\mathrm{res}}(t_f)$ & day 81, end of follow-up & Is the reservoir still detectable? \\
$R_b(t_L)$ and $R_b(t>t_L)$ & day 67 and follow-up & Did bacteremia clear and then return? \\
\bottomrule
\end{tabular}
\end{center}
Here $t_L$ is linezolid completion and $t_f$ is the end of the 14-day
follow-up. If one wished to describe only whether blood became detectable
during follow-up, a separate quantity would be
\begin{equation}
H(\theta)=\max_{t_L<t\leq t_f}R_b(t;\theta)-\LOD.
\label{eq:blood-followup}
\end{equation}
Then $H(\theta)\geq0$ means that bacteremia was detectable at least once
during follow-up. Even this is not sufficient by itself to define relapse:
relapse additionally requires $R_b(t_L)<\LOD$, so that blood was cleared when
linezolid ended before becoming detectable again.

\subsection{How the bracket is found}

A deterministic grid is evaluated first. Adjacent values $a<b$ are selected
when their classifications differ, equivalently when
\begin{equation}
F(a)F(b)<0.
\label{eq:sign-bracket}
\end{equation}
This establishes that the switch lies somewhere inside $[a,b]$, provided the
classification changes only once within that interval. The grid is therefore
used to establish direction, monotonicity, and a safe bracket; its midpoint
is not accepted as the final switch estimate.

\subsection{What Brent's algorithm does}

The bracket is passed to SciPy's
\texttt{scipy.optimize.brentq} implementation \cite{scipyBrentq}. Brent's
method combines three ideas \cite{brent1973}:
\begin{enumerate}
\item \textbf{Bisection} tries the midpoint of the bracket. It is slow but
reliable because it always reduces a valid sign-changing interval.
\item \textbf{Secant interpolation} uses the two endpoint values to estimate
where a zero would lie and is often faster than bisection.
\item \textbf{Inverse quadratic interpolation} uses three recent evaluations
to make an even faster estimate when their geometry is suitable.
\end{enumerate}
At each iteration, the algorithm accepts an interpolation step only when it
is judged safe; otherwise it falls back to bisection. After evaluating the
new point, it discards the half of the interval that does not contain the sign
change. Repetition contracts the bracket until its width satisfies the
requested absolute parameter tolerance.

The dedicated deterministic scripts use an absolute tolerance of $10^{-3}$
in the corresponding parameter units. The immune Monte Carlo script uses a
tighter $10^{-6}~\mathrm{h}^{-1}$ tolerance for its displayed switch. These
differences are smaller than the two-significant-figure precision used to
report the transitions.

\subsection{Important limitation: this is a discontinuous boundary}

The SciPy documentation specifies continuity as an assumption for
\texttt{brentq} \cite{scipyBrentq}. Equation~\eqref{eq:brent-objective} is
not continuous under the hard-LOD rule: terminal burden can jump from zero to
thousands of CFU/mL without ever taking the value 10~CFU/mL. Consequently,
the usual theorem that Brent's method converges to a root satisfying
$F(\theta)=0$ does not apply here.

In this analysis, \texttt{brentq} is used as a bracket-contraction procedure
that numerically localizes a binary clearance--persistence change. The result
should be called a \emph{classification boundary}, \emph{regime switch}, or
\emph{persistence threshold}; it should not be described as a conventional
continuous root.

Nor has the switch been demonstrated to be a transcritical bifurcation. A
transcritical bifurcation requires two equilibrium branches to meet and
exchange stability. Establishing that claim would require equilibrium
continuation and a Jacobian/eigenvalue stability analysis of the underlying
continuous ODE. The present calculation instead uses a finite treatment
course, a finite observation time, and discrete LOD events.

\subsection{Current deterministic switches}

\begin{table}[htbp]
\centering
\small
\begin{tabular}{@{}lllll@{}}
\toprule
Parameter & Grid & Points & Persistence side & Reported switch \\
\midrule
$\EC$ & $0.05$--$3.3~\mathrm{mg/L}$ & 131 & above & $1.3~\mathrm{mg/L}$ \\
$\EM$ & $0.70$--$2.2~\mathrm{h}^{-1}$ & 111 & below & $1.5~\mathrm{h}^{-1}$ \\
$\KI$ & $0.02$--$0.65~\mathrm{h}^{-1}$ & 127 & below & $0.047~\mathrm{h}^{-1}$ \\
\bottomrule
\end{tabular}
\caption{Current resistant-reservoir classification boundaries. The unrounded
computational values, rather than the displayed values, are used in subsequent
probability calculations.}
\label{tab:switches}
\end{table}

\section{Monte Carlo procedure}

\subsection{Meaning of a Monte Carlo iteration}

One iteration represents one virtual patient. The patient receives one random
value of the parameter under study, while the treatment protocol, starting
state, solver settings, and every other biological parameter remain fixed.
The model is then integrated for the entire 81-day course. Repeating this
procedure 1000 times produces a distribution of complete patient trajectories
and outcomes.

This is parameter-variability propagation: the random number does not add
noise to the ODE at each time step. Once a patient's parameter value is drawn,
that value remains constant for the patient's entire simulation.

\subsection{Why a mean-corrected log-normal distribution is used}

All three sampled parameters must be positive. They are drawn from log-normal
distributions. If $X$ has required arithmetic mean $m$ and coefficient of
variation $CV$, the code calculates
\begin{align}
\sigma_{\ln}&=\sqrt{\ln(1+CV^2)},\\
\mu_{\ln}&=\ln(m)-\frac{1}{2}\sigma_{\ln}^2,\\
X&\sim\operatorname{LogNormal}(\mu_{\ln},\sigma_{\ln}^2).
\label{eq:lognormal}
\end{align}
The subtraction of $\sigma_{\ln}^2/2$ is the mean correction. It ensures
\begin{equation}
\mathbb{E}[X]=\exp\!\left(\mu_{\ln}+\frac{\sigma_{\ln}^2}{2}\right)=m.
\end{equation}
Without this correction, entering $\ln(m)$ as the log-scale location would
make $m$ the median, not the mean, and the Monte Carlo calibration would no
longer match the fixed baseline.

\begin{table}[htbp]
\centering
\small
\begin{tabular}{@{}llll@{}}
\toprule
Analysis & Arithmetic mean & CV & Random seed \\
\midrule
$\EC$ & $1.13~\mathrm{mg/L}$ & $0.60$ & 44 \\
$\EM$ & $1.65~\mathrm{h}^{-1}$ & $0.080$ & 44 \\
$\KI$ & $0.10~\mathrm{h}^{-1}$ & $0.80$ & 42 \\
\bottomrule
\end{tabular}
\caption{Centralized Monte Carlo distributions. Each analysis contains 1000
draws and varies only the listed parameter.}
\label{tab:mc-distributions}
\end{table}

\subsection{Exact computational sequence}

Each Monte Carlo script follows the same essential sequence:
\begin{enumerate}
\item Load the model, centralized parameters, canonical time grid, hybrid
solver, and outcome-classification function.
\item Construct the vancomycin and linezolid concentration functions once.
\item Read the selected parameter's arithmetic mean, CV, and random seed from
\texttt{parameters.py}.
\item Convert the mean and CV to $\mu_{\ln}$ and $\sigma_{\ln}$ using
Eq.~\eqref{eq:lognormal}.
\item Initialize a dedicated NumPy random-number generator and draw all 1000
parameter values. Drawing them once makes the cohort reproducible.
\item Pre-allocate four arrays with 1000 rows and 1453 columns, one array for
each bacterial state. Row $i$ will contain virtual patient $i$'s trajectory.
\item For patient $i$, copy the fixed baseline and override only the sampled
parameter. For the immune analysis, construct a new
\texttt{ImmuneResponse} object containing the sampled $\KI$.
\item Integrate the complete four-state model over all 1453 reporting times.
The hybrid solver has already applied extinction and reactivation events.
\item Store all four trajectories, then repeat until all 1000 patients are
complete.
\item Apply the same trajectory-based outcome definitions to every patient's
resistant-blood history.
\item Summarize counts, percentages, distributions, and time courses, and
write the figures with parameter-provenance metadata.
\end{enumerate}

\subsection{Clinical outcome definitions}

Let $t_L$ denote linezolid completion and let the subsequent interval be the
14-day follow-up. The three outcomes are mutually exclusive:
\begin{align}
\text{relapse}:\quad &R_b(t_L)<\LOD
\quad\text{and}\quad \max_{t>t_L}R_b(t)\geq\LOD,\\
\text{persistent failure}:\quad &R_b(t_L)\geq\LOD,\\
\text{resolution}:\quad &R_b(t_L)<\LOD
\quad\text{and}\quad R_b(t)<\LOD\ \text{for every }t>t_L.
\end{align}
A detectable reservoir without a post-treatment return of bacteremia is not
called relapse. Conversely, a reservoir can cause relapse by remaining
detectable through treatment and reseeding blood after linezolid ends.

The switch threshold is not used to force an outcome label. Every label is
calculated from that patient's simulated blood trajectory. This distinction
is important near the switch and whenever blood remains detectable at the end
of treatment. Section~\ref{sec:sign-meaning} explains why the sign of the
reservoir-based Brent objective cannot be interpreted directly as the presence
or absence of post-linezolid bacteremia.

\subsection{Trajectory plots and missing-value display convention}

The model stores extinct populations as exact zeros. A logarithmic axis cannot
display zero, so plotting code replaces zeros with \texttt{NaN} only for
display and NaN-aware summaries. This does not change the simulation or the
outcome classification.

For the EC50 and E-max kinetics figures, the median and 5th--95th percentile
band are calculated within the patients who relapsed, making the reseeding
trajectory visible. The immune-response figure calculates NaN-aware summaries
over all patients; at each time these are conditional on the compartment being
detectable. These conditional summaries must not be interpreted as
unconditional cohort quantiles in which extinct patients have a small positive
burden.

\section{From switch distance to relapse probability}

\subsection{Why the raw distance is insufficient}

The simple difference $T-m$ is not comparable across parameters because the
parameters have different units and different CVs. A difference of 0.10 is
large for a narrow distribution and small for a broad distribution. Because
the sampling distributions are log-normal, distance is measured after taking
natural logarithms and is divided by the log-scale standard deviation.

For an unrounded deterministic switch $T$, define
\begin{equation}
z_T=\frac{\ln(T)-\mu_{\ln}}{\sigma_{\ln}}.
\label{eq:z}
\end{equation}
The z-value answers: ``How many log-scale standard deviations is the switch
from the center of the sampling distribution?'' The standard normal
cumulative distribution function
\begin{equation}
\Phi(z)=P(Z\leq z),\qquad Z\sim N(0,1),
\end{equation}
converts that standardized distance into an area under the distribution.

\subsection{Choosing the correct tail}

For $\EC$, high values weaken linezolid, so persistence is the upper tail:
\begin{equation}
P(\text{persistence side})=P(\EC>T)=1-\Phi(z_T).
\end{equation}
For $\EM$ and $\KI$, low values weaken killing or host clearance, so
persistence is the lower tail:
\begin{equation}
P(\text{persistence side})=P(X<T)=\Phi(z_T).
\end{equation}

For example, the immune switch has $z=-0.74$. Approximately 23.1\% of a
standard normal distribution lies to the left of $-0.74$. Therefore about
23.1\% of the configured immune distribution is expected to lie in the weak-
immunity persistence regime. The seeded finite cohort placed 240 of 1000
patients (24.0\%) on that side; a finite random sample need not equal the
theoretical percentage exactly.

\subsection{Current distance and outcome calculations}

\begin{table}[htbp]
\centering
\small
\begin{tabular}{@{}lrrrrr@{}}
\toprule
Parameter & $z_T$ & Theoretical tail & Sampled side & Relapse & Persistent \\
\midrule
$\EC$ & $0.60$ & 27.5\% & 300 (30.0\%) & 289 (28.9\%) & 11 (1.1\%) \\
$\EM$ & $-0.75$ & 22.6\% & 220 (22.0\%) & 220 (22.0\%) & 0 (0.0\%) \\
$\KI$ & $-0.74$ & 23.1\% & 240 (24.0\%) & 240 (24.0\%) & 0 (0.0\%) \\
\bottomrule
\end{tabular}
\caption{Relationship between the deterministic switch, probability mass on
the persistence side, and trajectory-defined clinical outcomes. The z-values
use unrounded computational switches.}
\label{tab:z-outcomes}
\end{table}

For each analysis in Table~\ref{tab:z-outcomes}, the number on the persistence
side equals relapse plus persistent treatment failure. Thus the switch-tail
area estimates the \emph{relapse-prone} fraction, not necessarily the relapse
rate itself. A persistence-side patient is counted as relapse only if blood is
below the LOD when linezolid ends and later becomes detectable. If blood is
already detectable when linezolid ends, the outcome is persistent treatment
failure.

This explains why persistence-side fractions and relapse percentages can
differ. In the immune analysis all 240 persistence-side patients, and in the
E-max analysis all 220, first cleared blood and then relapsed. In the EC50
analysis, 11 of the 300 persistence-side patients never cleared blood by
treatment completion.

\section{Analysis-specific interpretation and verified results}

\begin{table}[htbp]
\centering
\small
\begin{tabular}{@{}lrrr@{}}
\toprule
Analysis & Relapse & Persistent failure & Resolution \\
\midrule
$\EC$ & 289 (28.9\%) & 11 (1.1\%) & 700 (70.0\%) \\
$\EM$ & 220 (22.0\%) & 0 (0.0\%) & 780 (78.0\%) \\
$\KI$ & 240 (24.0\%) & 0 (0.0\%) & 760 (76.0\%) \\
\bottomrule
\end{tabular}
\caption{Verified mutually exclusive outcomes for 1000 virtual patients per
analysis. Each row sums to 1000.}
\label{tab:outcomes}
\end{table}

\paragraph{Linezolid EC50.}\mbox{}\par
Increasing $\EC$ reduces potency. The clearance--persistence switch is
$1.3~\mathrm{mg/L}$, and the sampling distribution has arithmetic mean
$1.13~\mathrm{mg/L}$ and CV 0.60. The high-EC50 upper tail is therefore the
persistence-prone group.

\paragraph{Linezolid E-max.}\mbox{}\par
Decreasing $\EM$ lowers the maximum killing capacity. The reported switch is
$1.5~\mathrm{h}^{-1}$, and the sampling distribution has arithmetic mean
$1.65~\mathrm{h}^{-1}$ and CV 0.080. Because this distribution is narrow, a
small raw separation between the mean and switch can still correspond to a
substantial standardized distance and tail probability.

\paragraph{Immune response.}\mbox{}\par
Decreasing $\KI$ weakens host clearance. The switch is
$0.047~\mathrm{h}^{-1}$, while the distribution mean is
$0.10~\mathrm{h}^{-1}$ with CV 0.80. This broad distribution places a
clinically important lower tail below the switch even though the arithmetic
mean is more than twice the reported switch.

All three relapse percentages fall within the prespecified 20--30\%
calibration range. They are outcomes of the chosen means, CVs, other fixed
parameters, hard-LOD logic, and relapse definition; they should not be
interpreted as independently validated clinical predictions.

\section{Reproducibility map}

\begin{longtable}{@{}p{0.28\textwidth}p{0.66\textwidth}@{}}
\toprule
File & Role \\
\midrule
\endhead
\path{parameters.py} & Centralized parameters, sampling distributions,
seeds, time grid, outcome definitions, and solver configuration. \\
\path{model_Bacteremia.py} & Pharmacokinetics, immune-response object, and
four-state ODE right-hand side. \\
\path{solver.py} & LSODA integration with hard-LOD extinction and
transfer-driven reactivation events. \\
\path{EC50_lzd_threshold_sweep.py} & Deterministic EC50 switch. \\
\path{Emax_lzd_threshold_sweep.py} & Deterministic E-max switch. \\
\path{ImmuneResponse_threshold_sweep.py} & Deterministic immune switch. \\
\path{MC_ec50_lzd.py} & EC50 Monte Carlo analysis and figures. \\
\path{MC_emax_lzd.py} & E-max Monte Carlo analysis and figures. \\
\path{MC_immune_response.py} & Immune-response Monte Carlo analysis and
figures. \\
\bottomrule
\end{longtable}

The principal figures are:
\begin{itemize}
\item \path{mc_ec50_lzd_sweep_kinetics.png},
      \path{mc_emax_lzd_sweep_kinetics.png}, and
      \path{mc_immune_response.png};
\item the three corresponding switch/population figures; and
\item \path{distribution_to_switch.png}, \path{stacked_outcomes.png}, and
      \path{relapse_time_curve.png}.
\end{itemize}

\section{Interpretive cautions}

\begin{enumerate}
\item The three Monte Carlo analyses are one-at-a-time analyses. They do not
represent correlations or simultaneous variability among EC50, E-max, and
immune response.
\item The CVs are outcome-calibrated inputs. They are not claimed to be
population estimates obtained from an independent clinical dataset.
\item A deterministic reservoir switch is mechanistic context, not the
clinical relapse definition.
\item The hard LOD is a modeling rule. It changes the mathematics from a
fully continuous ODE to a hybrid dynamical system.
\item Brent's continuity-based root guarantee does not apply to the
hard-LOD objective. The returned value localizes a classification change.
\item Calling the switch a transcritical bifurcation would require a separate
equilibrium and stability analysis that has not been performed.
\item Displayed parameters and thresholds are rounded, but probability
calculations use unrounded computational values. This is especially important
for the narrow E-max distribution.
\end{enumerate}

\bibliographystyle{ieeetr}
\bibliography{references}

\end{document}
